\documentclass[10pt]{resume}
\usepackage[left=0.5in,top=0.4in,right=0.5in,bottom=0.4in]{geometry}
\usepackage[colorlinks=true,urlcolor=blue]{hyperref}
\usepackage{enumitem}

\name{LOÏC ARON MBASSI EWOLO}
\address{Stage Ingénieur : R\&D Électronique \& Systèmes Embarqués (4 à 5 mois dès Avril 2026)}
\address{
Cergy, France \\
\href{mailto:loicaron.mbassiewolo@ensea.fr}{loicaron.mbassiewolo@ensea.fr} \,|\, 
\href{tel:+33759523398}{+33 7 59 52 33 98} \\
\href{https://www.linkedin.com/in/loic-mbassi/}{LinkedIn} \,|\, 
\href{https://github.com/Nameless0l}{GitHub}
}

\begin{document}

%--------------------------------------------------
% RÉSUMÉ PROFESSIONNEL / PROFIL
%--------------------------------------------------
\begin{rSection}{Résumé Professionnel}
Élève-ingénieur en double diplôme ENSEA/Polytechnique Yaoundé, spécialisé en systèmes embarqués et IA. Je recherche un stage R\&D chez VALOTECH pour appliquer mes compétences en électronique numérique, programmation bas niveau (C, VHDL) et soudure de composants.
\end{rSection}

%--------------------------------------------------
\begin{rSection}{Éducation}
{\bf ENSEA - École Nationale Supérieure de l'Électronique} \hfill {\em 2025--2027}\\
\textbf{Double diplôme d'ingénieur -- Systèmes Embarqués \& IA}\\
Cursus : FPGA/VHDL, Traitement du Signal, Deep Learning, Architecture processeurs.

{\bf École Nationale Supérieure Polytechnique de Yaoundé (ENSPY)} \hfill {\em 2021--2025}\\
\textbf{Diplôme d'ingénieur de conception en informatique} (Niveau Master 2)\\
\textit{Cursus :} Mathématiques Appliquées, Génie Logiciel, Systèmes temps réel.
\end{rSection}

%--------------------------------------------------
\begin{rSection}{Compétences Techniques}
\begin{tabular}{ @{} >{\bfseries}l @{\hspace{6ex}} l }
Embarqué / Hardware & C/C++, VHDL (FPGA), STM32, Soudure, Protocoles I2C/SPI, Linux \\
Maths / Algorithmes & SLAM, Monte Carlo, Algèbre Linéaire, Optimisation \\
IA / Data & PyTorch, TensorFlow, Computer Vision (YOLOv10), NLP \\
Outils & Git, Docker, Quartus, ModelSim, Altium/KiCad (PCB) \\
Langues & Français (Natif), Anglais (Professionnel/Technique)
\end{tabular}
\end{rSection}

%--------------------------------------------------
\begin{rSection}{Projets Académiques}

{\bf Harmony Gloves -- Traduction Langue des Signes} \hfill {\em 2024}\\
Prototypage de gants instrumentés au Laboratoire IOTA (Polytechnique). Contribution à la conception du circuit et \textbf{soudure} des composants (Capteurs IMU/Flex). Prise de données expérimentales pour l'entraînement du modèle TinyML.

{\bf Upgrade Jetson -- Challenge CoHoMa (Défense)} \hfill {\em En cours}\\
Programmation d'un robot autonome pour cartographie en environnement non structuré. Implémentation d'algorithmes de \textbf{SLAM} et méthodes de \textbf{Monte Carlo}. Traitement Lidar 3D VLP16 et caméra de profondeur, développement C/C++.

{\bf Conception FPGA (VHDL) -- ENSEA} \hfill {\em 2025}\\
Développement de descriptions matérielles en \textbf{VHDL}. Simulation et validation comportementale sur Quartus/ModelSim.

{\bf Projet Taxi -- Simulation Système Temps Réel} \hfill {\em 2024}\\
Simulation flotte taxis multiprocessus en C. Gestion mémoire partagée, IPC (Signaux), ordonnanceur FIFO. Visualisation temps réel avec OpenGL.

\end{rSection}

%--------------------------------------------------
\begin{rSection}{Expérience Professionnelle}
{\bf Assistant de Recherche -- MILA (Institut Québécois d'IA), Québec} \hfill {\em Juin -- Sept. 2025}\\
Recherche sur la généralisation tardive (Grokking) des réseaux de neurones. Implémentation de pipelines de tests statistiques sous \textbf{PyTorch}. Reproduction d'expériences SOTA.

{\bf Développeur Backend -- CSC Fintech, Yaoundé} \hfill {\em Oct. 2024 -- Jan. 2025}\\
Conception backend d'une plateforme de transactions sécurisées (Laravel, MySQL, Docker, Redis).
\end{rSection}

%--------------------------------------------------
\begin{rSection}{Distinctions et Engagements}

\textbf{1\textsuperscript{ère} Place -- CITS National Hackathon} \hfill \textit{2024} \\
Solution IA pour gestion intelligente des déchets urbains (optimisation routes -20\%).

\textbf{1\textsuperscript{ère} Place -- Mountain Hub Regional} \hfill \textit{2024} \\
Innovation technologique -- Projet sélectionné parmi 50+ équipes.

\textbf{Chef de Cellule Projets -- Polytechnique Yaoundé} \hfill \textit{2023--2025} \\
Management d'équipes techniques, animation workshops (Laravel, React, ML).
\end{rSection}

\end{document}
